\chapter{Asynchronous Models}

\section{Core assumption}
\begin{defbox}
An asynchronous system places no fixed upper bound on message delay, process speed, or clock drift. Correctness must not depend on timing guarantees.
\end{defbox}

\section{Typical properties}
\begin{itemize}
  \item Messages may be delayed arbitrarily long.
  \item A process may run much slower than its peers.
  \item Timeouts are only heuristics, not proof of failure.
\end{itemize}

\section{Why it matters}
\begin{keybox}
Asynchrony is the most conservative model, so protocols that work here are usually robust under timing uncertainty.
\end{keybox}

\section{Exam prompts}
\begin{itemize}
  \item Why can an asynchronous algorithm not rely on bounded delays?
  \item What is the practical impact of using timeouts in this model?
  \item Give an example of a problem that becomes harder in a fully asynchronous setting.
\end{itemize}
